\section{What the new source changes about the order question}
\label{sec:order-research}

The new cached coefficients improve a certified exponential constant.
They leave the worst-case scale
$n\log\log n/\sqrt{\log n}$ unchanged.  The user-supplied Anthropic
development introduces a different algebraic mechanism, and therefore
deserves more than a comparison of its headline bound.  This section
examines a possible composition, makes one of the source's suggested
extensions precise, and proves a restricted obstruction to a natural
implementation of that extension.

\subsection{A direct composition spends the apparent extra saving}

The earlier precision packet gives an exact rank encoding of the
structured mode sum \cite{precision}.  If the first deformation uses
$r$ groups, its existing collision-packing certificate allows a retained
level range of order
\[
 u\asymp\frac{\log M}{\log r}.
\]
Its compressed mode-index graph consequently has a certified bit budget
on the scale
\begin{equation}
 B\asymp\frac{M}{\sqrt u}
       \asymp\frac{M\sqrt{\log r}}{\sqrt{\log M}}.
 \label{eq:composition-budget}
\end{equation}
These are the scales of that proof's envelopes, not lower bounds on the
actual entropy of every instance.

At fixed or polylogarithmic $r$, a group tensor can have degree comparable
to $B/r$, so the original bounded-degree contraction theorem cannot
simply be invoked again at size $B$.  This missing interface is already
identified in the earlier reports \cite{precision,barriers}.
Anthropic's complete-graph algorithm does permit arbitrary integer-valued node tensors
with charged entry access \cite[Theorem A]{anthropic}.  It is therefore a
more appropriate possible backend, after balancing and representation
costs have been supplied.

However, the inspected backend has an exponent rate of order
$e(r)\asymp1/\sqrt{\log r}$ along its certified schedule.  Even the
optimistic calculation that leaves out balancing and entry-evaluation
overheads gives
\begin{equation}
 e(r)B\asymp
 \frac1{\sqrt{\log r}}
 \frac{M\sqrt{\log r}}{\sqrt{\log M}}
 =\frac{M}{\sqrt{\log M}}.
 \label{eq:composition-cancellation}
\end{equation}
For fixed $r$ the backend changes a constant.  Increasing $r$ spends
the corresponding saving in the first-stage retained-level budget.
Allowing $r$ to grow would also require uniform accounting for the
exact-weight arrays and the newly expensive group entries.  Equation
\eqref{eq:composition-cancellation} is thus a diagnosis of this proposed
composition, not a completed composition theorem or an algorithmic lower
bound.  The new source resolves an important admissibility mismatch but
does not automatically supply the missing $M/\log M$ summation method.

\subsection{A cost-complete form of the source's growth suggestion}

Appendix B of the Anthropic paper builds crossing decompositions from
operators on the graded space
$V=\mathbb Q[t_1,t_2,\ldots]$, with $\operatorname{weight}(t_j)=j$.
Here $\dim V_E=p(E)$ is the ordinary partition number.  Its two operator
orders pass through spaces of weights at most $2h+s$ and $2h-s$.
The asymptotic relation $\log_2p(E)=A\sqrt E+O(\log(E+2))$ gives the
rank-sum parameter $\rho=1/\sqrt2$ and the logarithmic exponent
$\log_2(1/\rho)=1/2$ \cite[Appendices A and B]{anthropic}.

The source itself observes in Appendix H that growth like $E^\beta$
would instead give $\rho=2^{\beta-1}$.  The following proposition
formalizes that suggestion with the resource hypotheses needed for an
algorithm.  The growth-to-rate idea belongs to the source.  No algebra
meeting the stronger hypotheses below is constructed in this report.

\begin{proposition}[Conditional graded-growth transfer]
\label{prop:graded-growth}
Fix a rational $0<\beta<1$.  Suppose an effectively presented graded rational
operator system $V=\bigoplus_{E\ge0}V_E$ has fixed bases, finite,
nondecreasing dimensions $d_E\ge1$, and
\begin{equation}
 \log_2d_E=A E^\beta+O(\log(E+2)),\qquad A>0.
 \label{eq:graded-growth}
\end{equation}
For each bit length $b$, let $h$ be the least integer with $d_h\ge2^b$
and label $2^b$ distinct basis states of $V_h$ by $x\in\{0,1\}^b$.
Assume the following additional properties, uniformly in $b$.
\begin{enumerate}
\item A nondegenerate dual pairing and an exact operator identity give
operators $Q_s(x,v)$ on $V_h$, for $|s|\le2h$, and a positive integer
$D_h$, such that
\[
 \sum_{s=-2h}^{2h}Q_s(x,v)=D_h[x=v]\operatorname{Id}_{V_h}.
\]
\item Each $Q_s(x,v)$ has two exact factorizations, respectively as a
sum of $h^{O(1)}$ products $B_{v,s,j}A_{x,s,j}$ and as a sum of
$h^{O(1)}$ products $\widetilde A_{x,s,j}\widetilde B_{v,s,j}$.
The intermediate spaces in the first factorization have weights at
most $2h+s$, and those in the second at most $2h-s$.
All factorization entries are integral, after any specified common
denominator clearing.  This exact two-order identity is a hypothesis;
a dimension estimate alone does not imply it.
\item There are uniform algorithms computing $h$, labels, index lists,
$D_h$, piece values and factorization entries in $2^{O(b)}$ bit
operations.  The number of pieces is $b^{O(1)}$, and all integers have
$b^{O(1)}$ bits.  In particular the hypotheses include the basis and
rank-decoding costs and effective exact locality, wherever locality is
used to prove the preceding identities.
\end{enumerate}
Then the system supplies pieces in Anthropic's Definition 7 with
$\rho=2^{\beta-1}$.  Relative to its Theorem A and dense reduction, this
gives, for fixed width $k$,
\[
 \poly(n,m)2^{O_k(n\log\log n/(\log n)^{1-\beta})}
\]
time at arbitrary density and
$\poly(n,m)2^{O_{k,D}(n/(\log n)^{1-\beta})}$ when $m\le Dn$.
No polynomial-space conclusion is part of this implication.
\end{proposition}

\begin{proof}
For a piece define
$A_s(x,y,u,v)=\langle e_u^*,Q_s(x,v)e_y\rangle$, where $e_u^*$ is
the dual basis vector.
The first identity gives
$\sum_s A_s=D_h[x=v][y=u]$.
The first factorization separates $(x,y)$ from $(u,v)$, and the
second separates $(x,u)$ from $(y,v)$.  By monotonicity of $d_E$,
their respective numbers of rank terms are at most
$h^{O(1)}d_{2h+s}$ and $h^{O(1)}d_{2h-s}$.

The matching upper and lower estimates in \eqref{eq:graded-growth},
together with minimality of $h$, give
\[
 h=\Theta(b^{1/\beta}),\qquad
 A h^\beta\le b+O(\log(b+2)).
\]
Indeed $d_{h-1}<2^b$ supplies the latter inequality up to
$A(h^\beta-(h-1)^\beta)=O(1)$.  Rounding the logarithms of the
rank term counts to bits and using concavity of $x^\beta$ gives
\begin{align*}
 c_s^{\rm row}+c_s^{\rm col}
 &\le A\bigl((2h+s)^\beta+(2h-s)^\beta\bigr)+O(\log h)\\
 &\le 2A(2h)^\beta+O(\log h)\\
 &\le 2^{\beta+1}b+O(\log(b+2)).
\end{align*}
Thus Definition 7 holds with $4\rho=2^{\beta+1}$; its other
conditions are precisely the effective hypotheses above.  Finitely
many small values of $b$ can be absorbed by enlarging its fixed
constant.  Theorem A has logarithmic exponent
$\log_2(1/\rho)=1-\beta$.  For fixed rational $\beta$, the rational
rate approximations required by that theorem are computable by integer
power comparisons, so its stated consequences apply.
\end{proof}

For example, $\beta=1/3$ would replace the square-root denominator by
$(\log n)^{2/3}$.  This is a sufficient-condition calculation, not an
achieved improvement.  Even $\log d_E=\Theta(E^\beta)$ is insufficient
to identify the displayed $\rho$: unrelated constants in the upper
and lower estimates would remain in the rank bound.

\subsection{Why deleting oscillator modes does not realize it}

A natural candidate retains only square-weight generators,
$\mathbb Q[t_1,t_4,t_9,\ldots]$.  Its attraction is the smaller
partition growth.  Its immediate defect is that the current's positive
commutator series has the form
\begin{equation}
 C(z)=\sum_{j\ge1}j^2z^{j^2},
 \label{eq:square-current}
\end{equation}
up to a monomial shift or nonzero coefficient normalization.
It cannot satisfy the polynomial-clearing locality used in Appendix B.
The obstruction extends beyond this one choice.

\begin{lemma}[Rational series cannot have unbounded support gaps]
\label{lem:rational-gaps}
Let $C(z)=\sum_{n\ge0}c_nz^n$ over a field be rational and not a
polynomial.  After a finite prefix, its nonzero coefficient support has
bounded gaps.
\end{lemma}
\begin{proof}
Write $Q(z)C(z)=P(z)$ with $Q(0)=q_0\ne0$ and $d=\deg Q$.
For $n>\deg P$ the coefficients satisfy
\[
 q_0c_n+q_1c_{n-1}+\cdots+q_dc_{n-d}=0.
\]
If $d$ consecutive coefficients in this tail were zero, induction
would force every later coefficient to vanish.  This would make $C$
a polynomial.  The case $d=0$ already makes $C$ a polynomial and is
excluded.  Therefore sufficiently late gaps have length less than $d$.
\end{proof}

The series \eqref{eq:square-current} has arbitrarily long zero runs
between consecutive squares and infinitely many nonzero terms, so it is
not rational.  In particular, no nonzero polynomial $Q$ makes $QC$ a
polynomial.  In the independent-current construction, clearing the
two-variable commutator by a fixed power of $(\omega_1-\omega_2)$
would, on taking its positive expansion and setting
$z=\omega_2/\omega_1$, imply just such rationality.  More explicitly,
up to a monomial shift the two-sided scalar distribution is
$C(z)-C(z^{-1})$.  Multiplying its negative part by $(1-z)^N$ can
produce positive exponents only below $N$.  Thus annihilation of the
distribution forces all sufficiently high coefficients of
$(1-z)^NC(z)$ to vanish.  The square-weight
replacement therefore lacks the exact identity on which the claimed
rank bounds would depend.

\begin{proposition}[A growth obstruction for independent oscillators]
\label{prop:oscillator-growth}
Let $S\subseteq\mathbb Z_{>0}$ be infinite and be the nonzero
coefficient support, apart from finitely many indices, of a rational
nonpolynomial current series.  Give the polynomial space
\[
 V=\mathbb Q[t_n:n\in S]
\]
the grading $\operatorname{weight}(t_n)=n$.  There is $c>0$ such that,
for all sufficiently large $E$,
\begin{equation}
 \sum_{j=0}^E\dim V_j\ge2^{c\sqrt E}.
 \label{eq:cumulative-growth}
\end{equation}
Consequently this model cannot satisfy a uniform exact-degree upper
bound $\dim V_E\le\exp(O(E^\beta+\log(E+2)))$ with $\beta<1/2$.
\end{proposition}
\begin{proof}
By Lemma~\ref{lem:rational-gaps}, a fixed integer $d$ bounds the gaps
in $S$ after a finite prefix.  For sufficiently large $B$, at least
$B/(2d)$ members of $S$ lie in $[1,B]$.  Independently choosing
exponent zero or one on these generators produces at least
$2^{B/(2d)}$ distinct monomials, every one of weight at most
$B|S\cap[1,B]|\le B^2$.  Take $B=\lfloor\sqrt E\rfloor$.
This proves the cumulative, rather than pointwise, lower bound
\eqref{eq:cumulative-growth}.  On the other hand, the proposed
exact-degree upper bound would imply
\[
 \sum_{j=0}^E\dim V_j
 \le (E+1)\exp\bigl(O(E^\beta+\log(E+2))\bigr),
\]
contradicting \eqref{eq:cumulative-growth} when $\beta<1/2$.
\end{proof}

The independence in this proposition is not removed by taking a
quotient that preserves a nondegenerate Heisenberg vacuum
representation.  Suppose a nonzero vacuum $v_0$ is annihilated by
positive operators $a_n$, the negative operators $a_{-n}$ raise weight
by $n$, operators of either sign commute among themselves, and
\[
 [a_n,a_{-m}]=\delta_{nm}\gamma_n\operatorname{Id},
 \qquad \gamma_n\ne0\quad(n\in S)
\]
over a characteristic-zero field.  In a relation among creation
monomials of the same weight, apply
$\prod_na_n^{\alpha_n}$ for any exponent tuple $\alpha$ appearing
there.  Every different same-weight tuple has some smaller coordinate
and is killed.  The matching tuple leaves
$\prod_n\alpha_n!\gamma_n^{\alpha_n}\,v_0\ne0$.
Hence its coefficient in the relation is zero, and all creation
monomials are independent.  Such a quotient retains the growth
obstruction.

These statements concern the specified independent-current or
Heisenberg-vacuum model.  They do not rule out different noncentral
operator algebras, other exact recovery identities, compressed tensor
representations, or faster SAT algorithms.  In particular, they are
not general lower bounds for vertex algebras or for SAT.  Their useful
conclusion is narrower: merely deleting oscillator modes, or imposing
relations while retaining the displayed Heisenberg premises, does not
implement the source's smaller-growth suggestion.

\subsection{Further research questions}

\begin{enumerate}
\item Can a noncentral graded operator algebra supply the pairing,
two-order identity and effective factorization of
Proposition~\ref{prop:graded-growth} with $\beta<1/2$?  A candidate
must escape the preceding Heisenberg premises while preserving exact
recovery and uniform bit costs.
\item Can a different exact recovery method use a nonrational current?
Truncating its commutator is insufficient unless the residual is proved
to vanish in the actual counting ring and the required precision is
fully charged.
\item Can the mode-rank and clique algorithms be composed with complete
balancing and entry-query accounting to improve a coefficient or the
memory tradeoff?  An order claim must additionally overcome
\eqref{eq:composition-cancellation}.
\item Which relations among the \emph{actual} group tensors permit
implicit summation without materializing their formal boundary tables?
Existing inventory barriers concern a representation; they do not
establish independence or noncancellation of the values in those tables
\cite{precision,barriers}.
\end{enumerate}

The first two questions seek a new algebraic ingredient.  The latter
two seek a representation that uses more than pair entropy.  None of
these required interfaces has been supplied here, so the proved
positive improvement remains the smaller cached contraction constant
and the exact kernel refinement established elsewhere in this report.
